\documentclass[pdflatex,sn-mathphys,Numbered]{sn-jnl}%

\usepackage{graphicx}%
\usepackage{multirow}%
\usepackage{amsmath,amssymb,amsfonts}%
\usepackage{amsthm}%
\usepackage{mathrsfs}%
\usepackage[title]{appendix}%
\usepackage{xcolor}%
\usepackage{textcomp}%
\usepackage{manyfoot}%
\usepackage{booktabs}%
\usepackage{algorithm}%
\usepackage{algorithmicx}%
\usepackage{algpseudocode}%
\usepackage{listings}%
\usepackage{tikz}

\raggedbottom

\begin{document}

\title[ConstrumindNet: Spatio-Temporal Equipment Allocation]{ConstrumindNet: A Hybrid Spatio-Temporal Deep Learning and Latent-Refined Framework for Multi-Horizon Heavy Equipment Allocation and Fleet Scheduling}

\author[1]{\fnm{Danthuluri} \sur{Sudha}}
\author[1]{\fnm{Santosh Kumar} \sur{Jankatti}}
\author[1]{\fnm{Ramandeep} \sur{Kaur}}
\author*[2]{\fnm{Bondu} \sur{Venkateswarlu}}\email{bonduvenkat-cse@dsu.edu.in}
\author[2]{\fnm{T.} \sur{Gayathri}}
\author[3]{\fnm{G.} \sur{Manjula}}
\author[4]{\fnm{Harish H.} \sur{Kenchannavar}}

\affil[1]{\orgdiv{Department of Computer Science and Technology}, \orgname{Dayananda Sagar University}, \orgaddress{\city{Bangalore}, \state{Karnataka}, \country{India}}}
\affil*[2]{\orgdiv{Department of Computer Science and Engineering}, \orgname{Dayananda Sagar University}, \orgaddress{\city{Bangalore}, \state{Karnataka}, \country{India}}}
\affil[3]{\orgdiv{Department of Computer Science and Design}, \orgname{Dayananda Sagar Academy of Technology and Management}, \orgaddress{\city{Bangalore}, \state{Karnataka}, \country{India}}}
\affil[4]{\orgdiv{Department of Computer Science and Engineering (Data Science)}, \orgname{Dayananda Sagar College of Engineering}, \orgaddress{\city{Bangalore}, \state{Karnataka}, \country{India}}}

\abstract{Heavy equipment allocation and fleet scheduling across geographically distributed, multi-phase civil and industrial infrastructure projects represent critical operational challenges in modern construction management. Suboptimal machinery dispatching, unscheduled maintenance disruptions, idle transit bottlenecks, and cascading project delays generate immense economic losses and severe carbon emissions. Classical operational research techniques, including Mixed-Integer Linear Programming (MILP) and heuristic dispatching rules, suffer from exponential computational complexity under dynamic fleet constraints, fail to model non-Euclidean spatial mobilization networks, and lack the capacity to capture temporal dependencies across multi-stage construction phases. While recent deep learning paradigms (e.g., CNN-LSTM, spatio-temporal Graph Neural Networks, and Transformers) have advanced predictive scheduling, they remain prone to cumulative multi-horizon rollout drift, spatial allocation inconsistency, and high-dimensional constraint violations. To resolve these pressing bottlenecks, we introduce \textbf{ConstrumindNet}, a novel hybrid spatio-temporal deep learning framework augmented with Latent Feature Refinement (LFR) and operational constraint regularisation for multi-horizon equipment allocation and fleet scheduling. ConstrumindNet integrates: (1) a multi-scale temporal learning hierarchy decomposing scheduling dynamics into short-term tactical shifts (1--3 days), medium-term phase handovers (7--14 days), and long-term project horizons (30+ days); (2) a spatial graph attention network explicitly capturing regional mobilization cost-distance topologies; (3) a Latent Feature Refinement (LFR) module that conducts nonlinear feature mixing across compressed latent allocations to capture non-local project-machinery synergies; and (4) an operational-aware physics and consistency regularizer penalizing maintenance-schedule collisions, excessive transit mileage, and cross-site resource starvation. Empirical benchmarking across a synchronized multi-site operational dataset spanning 46 heavy industrial machines, 8 active project sites, and 5 sequential construction phases demonstrates that ConstrumindNet achieves an Allocation Match Accuracy (AMA) of 91.4\% at Day 1 and 84.2\% at Day 7, reducing multi-horizon Root Mean Squared Error (RMSE) by up to 21.3\% and cutting allocation conflict rates by 68.4\% without requiring intractable combinatorial solvers. Extensive ablation studies, cross-project regime transfer evaluations, and paired statistical significance tests ($p < 0.005$, Cohen's $d = 0.92$) confirm the robustness and real-time operational viability of the proposed architecture.}

\keywords{Heavy equipment allocation, Spatio-temporal deep learning, Fleet scheduling, Graph neural networks, Latent feature refinement, Operational regularisation, Multi-horizon optimization}

\maketitle

\section{Introduction}\label{sec1}
Capabilities to dynamically schedule and allocate heavy industrial machinery with high precision have become indispensable for large-scale civil infrastructure, commercial construction, and regional industrial projects \cite{hossain2025,russo2025}. Heavy equipment assets---including hydraulic excavators, dump trucks, heavy mobile cranes, high-capacity tower cranes, concrete boom pumps, transit mixers, motor graders, bulldozers, and deep piling rigs---represent immense capital investments whose underutilization or misallocation leads to severe project cost overruns, idle engine emissions, and cascading contractual delays \cite{chiranjeevi2025,inam2024}. In distributed multi-site construction networks, machinery must be continuously coordinated across multiple concurrent projects operating at disparate geographical locations and sequential construction phases (ranging from earthworks and deep foundation piling to superstructure erection, facade installation, and MEP finishing) \cite{inam2025a,inam2025b}.

Older operational research techniques, such as Mixed-Integer Linear Programming (MILP), branch-and-bound algorithms, and priority dispatching heuristics (e.g., Earliest Due Date, Shortest Processing Time), remain widely referenced for their theoretical transparency and mathematical optimality \cite{hossain2025,meng2025}. Nonetheless, these formulations rely on static linear assumptions and suffer from combinatorial explosion ($NP$-hard complexity) when applied to enterprise fleet scales under stochastic site disruptions \cite{wang2025,ejiyi2025}. Over the past few years, deep learning architectures, including convolutional neural networks (CNNs) \cite{russo2025}, recurrent neural networks (RNNs) \cite{chiranjeevi2025}, and self-attention Transformers \cite{inam2025a}, have shown promising capabilities in capturing non-linear temporal dynamics and equipment utilization patterns. However, purely data-driven sequence models remain prone to severe multi-horizon error propagation during recursive dispatch rollouts, fail to model non-Euclidean road logistics networks, and frequently produce physically infeasible schedules that violate strict maintenance windows or machinery capacity constraints \cite{hong2024,park2024}.

In an effort to overcome these drawbacks, recent studies have investigated hybrid learning architectures that combine structured domain knowledge with deep representation networks \cite{meng2025,wang2025}. In parallel, Graph Neural Networks (GNNs) have demonstrated marked advantages in capturing spatial interdependencies across distributed transportation networks \cite{hussain2023,shin2023}. Concurrently, latent-space feature refinement has emerged as a promising avenue for advancing predictive modeling, offering compact yet highly expressive nonlinear transformations in compressed embedding spaces \cite{li2025,liu2025}. Nevertheless, existing works remain fragmented: deep learning models are evaluated on isolated time-series tasks, while spatial network routing and hard operational constraint enforcement are handled separately through ad-hoc post-processing. A unified framework combining multi-scale spatio-temporal learning, graph-based teleconnection modeling, latent feature refinement, and operational constraint regularisation has yet to be explored.

Inspired by these challenges, this work introduces \textbf{ConstrumindNet}, a unified hybrid spatio-temporal deep learning framework for multi-horizon heavy equipment allocation and fleet scheduling. ConstrumindNet models the multi-site construction network as a dynamic spatio-temporal graph, utilizing multi-scale temporal decomposition to capture daily tactical shifts, medium-term phase handovers, and long-term project lifecycles. A spatial Graph Attention Network (GAT) models road logistics distances and regional site-depot teleconnections, while a Latent Feature Refinement (LFR) module enhances nonlinear feature mixing across compressed latent representations. Furthermore, operational-aware and regime-consistency regularisation terms are integrated to penalize maintenance schedule collisions, reduce transit churn, and ensure robust cross-project generalization.

The major contributions of this paper are:
\begin{enumerate}
\item A unified hybrid spatio-temporal deep learning architecture tailored for multi-horizon heavy equipment allocation and fleet scheduling across distributed infrastructure sites.
\item A multi-scale temporal learning hierarchy that simultaneously captures daily tactical shifts, construction phase handovers, and long-term project lifecycles.
\item A graph-based spatial attention mechanism modeling non-Euclidean road logistics topologies and equipment-sharing teleconnections.
\item A Latent Feature Refinement (LFR) module that induces expressive nonlinear feature mixing across compressed latent tokens while preserving computational scalability.
\item An operational-aware regularisation strategy enforcing preventive maintenance constraints, minimizing inter-site mobilization churn, and aligning latent distributions across heterogeneous project regimes.
\end{enumerate}

The remainder of this article is organized as follows. Section \ref{sec2} surveys related literature and presents a detailed comparative gap analysis. Section \ref{sec3} establishes the theoretical preliminaries and mathematical formulations. Section \ref{sec4} details the ConstrumindNet methodology, architectural components, multi-scale rollout workflow, operational regularisation, and formal algorithms. Section \ref{sec5} reports comprehensive experimental results, ablation studies, and statistical significance evaluations. Section \ref{sec6} discusses practical implications and study limitations, and Section \ref{sec7} concludes the paper with future research directions.

\section{Related Work}\label{sec2}
This section reviews prior literature across classical scheduling, machine learning, deep spatio-temporal networks, and constraint-aware learning.

\subsection{Classical Operations Research and Mathematical Programming}
Classical equipment dispatching models formulate resource allocation as variants of the Resource-Constrained Project Scheduling Problem (RCPSP) or Vehicle Routing Problem with Time Windows (VRPTW) \cite{hossain2025,meng2025}. While MILP solvers guarantee global optimality under deterministic conditions, their exponential computational complexity renders them intractable for real-time dispatching across large fleets and multi-month horizons \cite{wang2025}. Priority dispatching heuristics (e.g., Earliest Due Date) provide sub-second tractability but suffer from myopic local decisions that trigger downstream bottlenecks \cite{ejiyi2025}.

\subsection{Deep Sequential and Spatio-Temporal Models}
Data-driven deep learning models have been applied to industrial logistics and demand forecasting \cite{russo2025,chiranjeevi2025}. Recurrent architectures (LSTM, GRU) capture temporal dependencies in equipment operating hours \cite{inam2024}, while ConvLSTM models combine spatial convolutions with recurrent state transitions \cite{shin2023}. Recently, Transformer architectures with multi-head self-attention have achieved strong sequence forecasting performance \cite{inam2025a,li2025}. However, vanilla attention incurs quadratic complexity with token counts and exhibits cumulative error propagation when rolled out autoregressively over extended horizons \cite{liu2025}.

\subsection{Graph Neural Networks for Spatial Logistics}
Distributed construction networks exhibit non-Euclidean topologies where project sites and central depots are interconnected by highway networks \cite{hussain2023}. Graph Convolutional Networks (GCN) and Graph Attention Networks (GAT) allow explicit modeling of spatial teleconnections, road travel times, and mobilization logistics \cite{liu2025,yadav2025}. Incorporating graph attention enables neural models to reason over spatial proximity and regional fleet availability simultaneously \cite{tao2025}.

\subsection{Constraint-Aware and Domain-Informed Machine Learning}
Purely data-driven neural networks optimize statistical losses without awareness of physical boundaries, frequently generating schedules that violate maintenance windows or site parking limits \cite{meng2025,cuomo2022}. Physics-Informed Neural Networks (PINNs) address this by embedding differential governing laws into the training objective \cite{cuomo2022}. In industrial scheduling, operational constraint regularisation acts as a soft surrogate, penalizing constraint violations while preserving gradient-based trainability \cite{wang2025}.

A systematic summary of representative studies and identified research gaps is presented in Table \ref{tab1}.

\begin{table*}[t]
\caption{Summary of key related works and identified research gaps}\label{tab1}
\begin{tabular*}{\textwidth}{@{\extracolsep\fill}llllll@{\extracolsep\fill}}
\toprule
Study / Reference Focus & Core Methodology & Application Domain & Key Strength & Limitation & Research Gap \\
\midrule
Classical OR \cite{hossain2025,meng2025} & MILP, Branch-and-Cut & Project scheduling (RCPSP) & Exact mathematical optimality & Exponential complexity ($NP$-hard) & Dynamic, polynomial-time algorithms \\
Heuristics \cite{ejiyi2025,nguyen2024} & GA, PSO, Earliest Due Date & Fleet dispatching & Computationally tractable & Trapped in local optima & Real-time adaptive re-allocation \\
Deep Sequential \cite{russo2025,inam2024} & CNN, LSTM, CNN-LSTM & Telematics forecasting & Extracts non-linear temporal trends & Error propagation in multi-step rollouts & Drift-resilient sequence learning \\
Spatio-Temporal \cite{shin2023,li2025} & ConvLSTM, Transformers & Traffic \& logistics & Captures joint spatio-temporal features & Quadratic complexity; Euclidean grids & Non-Euclidean spatial graph modeling \\
Graph Learning \cite{hussain2023,liu2025} & GCN, GAT, Message Passing & Distributed networks & Preserves structural road topology & Lacks multi-scale temporal hierarchy & Unified graph-temporal frameworks \\
Constraint-Aware ML \cite{meng2025,cuomo2022} & PINNs, Soft Penalties & Dynamical systems & Enforces physical plausibility & Gradient imbalances during training & Multi-task operational regularizers \\
\textbf{ConstrumindNet (Proposed)} & \textbf{Multi-Scale GAT + LFR} & \textbf{Heavy equipment allocation} & \textbf{Joint graph, multi-scale \& constraints} & \textbf{Moderate architectural complexity} & \textbf{End-to-end drift-resilient scheduling} \\
\bottomrule
\end{tabular*}
\end{table*}

\section{Preliminaries}\label{sec3}
This section formalizes the spatio-temporal state tensor, spatial graph representation, and multi-objective optimization objectives.

\subsection{Spatio-Temporal Equipment Allocation Foundations}
Let $\mathcal{S} = \{s_1, \dots, s_{N_s}\}$ denote the set of $N_s$ discrete spatial locations ($N_p$ active project sites and $N_d$ regional depots). Let $\mathcal{E} = \{e_1, \dots, e_{N_e}\}$ represent the fleet of $N_e$ heavy machinery units across 12 equipment classes. At each discrete time step $t$, the global operational state is described by the tensor field:
\begin{equation}
X_t \in \mathbb{R}^{N_s \times N_e \times C} \label{eq1}
\end{equation}
where $C$ denotes the operational feature channel dimension. The forecasting and allocation task maps an observed historical context sequence of length $L$ to future states over horizon $H$:
\begin{equation}
\{X_{t-L+1}, \dots, X_t\} \xrightarrow{f_\Theta} \{\widehat{X}_{t+1}, \dots, \widehat{X}_{t+H}\} \label{eq2}
\end{equation}

\subsection{Spatial Logistics Graph and Latent Embeddings}
The multi-site logistics network is formalized as a weighted graph $\mathcal{G} = (\mathcal{V}, \mathcal{E}_g, \mathcal{W})$, where edge weights encode road mobilization distances and transit durations. Latent embeddings are extracted via patch tokenization and projected into embedding space:
\begin{equation}
Z \in \mathbb{R}^{N \times d} \label{eq3}
\end{equation}
To achieve expressive feature mixing without memory explosion, $Z$ is compressed into a compact bottleneck representation $Z_c \in \mathbb{R}^{N \times q}$ ($q \ll d$), refined through nonlinear transformation blocks, and residually fused:
\begin{equation}
Z^* = Z + \alpha Z_r \label{eq4}
\end{equation}
where $\alpha$ is a trainable residual scaling factor. Table \ref{tab2} summarizes the mathematical notations used throughout this paper.

\begin{table}[h]
\caption{Notation summary for ConstrumindNet}\label{tab2}
\begin{tabular}{@{}ll@{}}
\toprule
Symbol & Description \\
\midrule
$X_t \in \mathbb{R}^{N_s \times N_e \times C}$ & Global spatio-temporal fleet state tensor at time $t$ \\
$N_s, N_e, C$ & Spatial nodes ($N_s=7$), fleet units ($N_e=46$), feature channels ($C=32$) \\
$L, H$ & Input context window ($L=7$) and forecast horizon ($H=7$ to $30$) \\
$p, N$ & Spatial-fleet patch size ($p=8$) and token count ($N=648$) \\
$d, q$ & Latent embedding dimension ($d=128$) and bottleneck dimension ($q=16$) \\
$Z, Z_c, Z_r, Z^*$ & Encoder latent, compressed, refined, and fused representations \\
$\Theta_e, \Theta_r, \Theta_d$ & Trainable parameters of encoder, refinement module, and decoder \\
$\widehat{X}_{t+h}$ & Predicted equipment allocation state field at lead time $h$ \\
$\mathcal{L}_{\text{task}}, \mathcal{L}_{\text{oper}}, \mathcal{L}_{\text{cons}}$ & Task loss, operational constraint loss, and regime-consistency loss \\
$\lambda_1, \lambda_2$ & Regularization hyperparameters ($\lambda_1 = 0.15, \lambda_2 = 0.10$) \\
$\Delta_{\text{transfer}}$ & Cross-project regime transfer degradation metric \\
\bottomrule
\end{tabular}
\end{table}

\section{Materials and Methods}\label{sec4}
This section details the operational dataset schema, architectural modules, multi-scale iterative workflow, operational regularisation terms, and algorithmic representations.

\subsection{Problem Formulation and Objectives}
ConstrumindNet predicts multi-horizon states $\widehat{X}_{t+1:t+H}$ directly via point prediction. The primary task objective minimizes a priority-weighted Mean Squared Error:
\begin{equation}
\mathcal{L}_{\text{task}} = \frac{1}{H} \sum_{h=1}^H \sum_{i=1}^{N_s} \sum_{j=1}^{N_e} \omega_{i, j} \left\| X_{t+h}(i, j, :) - \widehat{X}_{t+h}(i, j, :) \right\|_2^2 \label{eq5}
\end{equation}
where $\omega_{i, j}$ prioritizes critical-path infrastructure projects. Multi-step recursive rollout is formulated as:
\begin{equation}
\widehat{X}_{t+k} = f_\Theta\left(\widehat{X}_{t+k-L:t+k-1}\right) \label{eq6}
\end{equation}

\subsection{Operational Datasets and Data Preprocessing}
The benchmark dataset captures synchronized operations across the Bangalore infrastructure corridor, comprising 46 heavy industrial machines, 8 major construction projects, and 7 spatial hubs. Table \ref{tab3} summarizes the dataset configuration.

\begin{table*}[t]
\caption{Operational datasets and preprocessing configuration for ConstrumindNet}\label{tab3}
\begin{tabular*}{\textwidth}{@{\extracolsep\fill}lllll@{\extracolsep\fill}}
\toprule
Dataset Name & Primary Operational Role & Sampling Timestep & Typical Telemetry Variables & Spatial Coverage \\
\midrule
Fleet Telematics & Real-time machine status & 1-hour / Daily stats & Operating hours, capacity, fuel rate, status & Central Depot \& 6 active sites \\
Project Schedules & Project priorities \& phases & Daily operational logs & Phase IDs, required equipment, delay notes & Whitefield, Airport Site, Hebbal \\
Maintenance CMMS & Downtime \& repair tracking & Event-driven logs & Overhaul type, start/end dates, workshop status & Central Maintenance Depot \\
GIS Road Network & Mobilization cost matrices & Static graph + dynamic & Pairwise road distances (km), mobilization hours & Bangalore Urban \& Peripheral Ring \\
\bottomrule
\end{tabular*}
\end{table*}

\subsection{ConstrumindNet Architecture Design}
The architecture integrates:
\begin{enumerate}
\item \textbf{Spatio-Temporal Graph Transformer Encoder:} Employs Multi-Head Self-Attention along temporal slices and Graph Attention (GAT) across spatial nodes:
\begin{equation}
A_{\text{spatial}}(i, j) = \frac{\exp\left(\text{LeakyReLU}\left(a^T [W h_i \,\|\, W h_j]\right)\right)}{\sum_{k \in \mathcal{N}_i} \exp\left(\text{LeakyReLU}\left(a^T [W h_i \,\|\, W h_k]\right)\right)} \label{eq7}
\end{equation}
where edges are sparsified with threshold $\tau = 0.30$.
\item \textbf{Latent Feature Refinement (LFR) Module:} Compresses latent tokens to dimension $q = 16$, executes multi-stage nonlinear feature-mixing transformations, and projects back to dimension $d = 128$ with residual gating (Eq. \ref{eq4}).
\item \textbf{Multi-Horizon Forecast Decoder:} Produces direct multi-step point predictions across all lead times $h \in \{1, \dots, H\}$ simultaneously:
\begin{equation}
\widehat{X}_{t+1:t+H} = g_\Phi(Z^*) \label{eq8}
\end{equation}
\end{enumerate}

\subsection{Operational-Aware Regularisation}
To enforce physical and contractual boundaries, we implement:
\begin{enumerate}
\item \textbf{Maintenance Collision Penalty:}
\begin{equation}
\mathcal{L}_{\text{maint}} = \frac{1}{H} \sum_{h=1}^H \sum_{i \in \mathcal{S}_{\text{projects}}} \sum_{j=1}^{N_e} M_{t+h}(j) \cdot \widehat{Y}_{t+h}(i, j) \label{eq9}
\end{equation}
\item \textbf{Mobilization Transit Penalty:}
\begin{equation}
\mathcal{L}_{\text{transit}} = \frac{1}{H-1} \sum_{h=1}^{H-1} \sum_{j, i, i'} \widehat{Y}_{t+h}(i, j) \widehat{Y}_{t+h+1}(i', j) D_{\text{transit}}(i, i') \label{eq10}
\end{equation}
\item \textbf{Cross-Regime Alignment Loss:}
\begin{equation}
\mathcal{L}_{\text{regime}} = \left\| \mathbb{E}_{\mathcal{D}_{\text{src}}}[Z^*] - \mathbb{E}_{\mathcal{D}_{\text{tgt}}}[Z^*] \right\|_2^2 \label{eq11}
\end{equation}
\end{enumerate}
The composite training loss is:
\begin{equation}
\mathcal{L}_{\text{total}} = \mathcal{L}_{\text{task}} + \lambda_1 \left(\mathcal{L}_{\text{maint}} + \beta \mathcal{L}_{\text{transit}}\right) + \lambda_2 \mathcal{L}_{\text{regime}} \label{eq12}
\end{equation}

\subsection{Algorithmic Procedures}
The computational procedures are formalized in Algorithms \ref{alg1}, \ref{alg2}, and \ref{alg3}.

\begin{algorithm}[h]
\caption{End-to-End Training Procedure for ConstrumindNet}\label{alg1}
\begin{algorithmic}[1]
\Require Dataset $\mathcal{D}$, context length $L$, horizon $H$, weights $\lambda_1, \lambda_2$.
\Ensure Optimized parameter set $\Theta^*$.
\State Initialize parameters $\Theta = \{\Theta_e, \Theta_r, \Theta_d\}$.
\State Construct spatial graph $\mathcal{G} = (\mathcal{V}, \mathcal{E}_g, \mathcal{W})$ from GIS distance matrix.
\For{epoch $e = 1$ to $E$}
    \For{each mini-batch $b = 1$ to $B$}
        \State Sample sliding windows $\{X_{t-L+1:t}\}$ and targets $\{X_{t+1:t+H}\}$.
        \State Compute latent embeddings $Z = \text{Encoder}(\text{PatchEmbed}(X), \mathcal{G})$.
        \State Compress tokens $Z_c = \text{Compress}(Z)$ and refine $Z_r = \text{Refine}(Z_c)$.
        \State Fuse residually: $Z^* = Z + \alpha Z_r$.
        \State Decode multi-horizon output $\widehat{X}_{t+1:t+H} = \text{Decoder}(Z^*)$.
        \State Compute total composite loss $\mathcal{L}_{\text{total}}$ via Eq. (\ref{eq12}).
        \State Update parameters $\Theta \leftarrow \Theta - \eta_t \nabla_\Theta \mathcal{L}_{\text{total}}$ via AdamW.
    \EndFor
    \State Evaluate validation loss $\mathcal{L}_{\text{val}}$; execute early stopping check.
\EndFor
\State \Return $\Theta^*$
\end{algorithmic}
\end{algorithm}

\begin{algorithm}[h]
\caption{Multi-Scale Iterative Forecasting with Latent Refinement}\label{alg2}
\begin{algorithmic}[1]
\Require Context sequence $\{X_{t-L+1:t}\}$, horizon $H$, scale set $\mathcal{S}$, model $\Theta^*$.
\Ensure Multi-horizon schedule $\widehat{X}_{t+1:t+H}$ and conflict alerts.
\For{each temporal scale $s \in \mathcal{S}$}
    \State Initialize rollout buffer $\mathcal{B}_0^{(s)} = \{X_{t-L+1:t}^{(s)}\}$.
    \For{step $k = 1$ to $H$}
        \State Compute latent embedding $Z_k^{(s)}$ and refined token $Z_{r, k}^{(s)}$.
        \State Decode single-step prediction $\widehat{X}_{t+k}^{(s)} = \text{Decoder}(Z_k^{(s)} + \alpha Z_{r, k}^{(s)})$.
        \State Append to buffer: $\mathcal{B}_k^{(s)} \leftarrow \mathcal{B}_{k-1}^{(s)} \cup \{\widehat{X}_{t+k}^{(s)}\}$.
    \EndFor
\EndFor
\State Synchronize and fuse multi-scale forecasts: $\widehat{X}_{t+h} = \sum_{s \in \mathcal{S}} \beta_s \widehat{X}_{t+h}^{(s)}$.
\State Audit predicted schedules against CMMS maintenance work orders.
\State \Return $\widehat{X}_{t+1:t+H}$ and Conflict Audit Report.
\end{algorithmic}
\end{algorithm}

\begin{algorithm}[h]
\caption{Formalized Cross-Project Regime Transfer Protocol}\label{alg3}
\begin{algorithmic}[1]
\Require Source dataset $\mathcal{D}_{\text{src}}$, target dataset $\mathcal{D}_{\text{tgt}}$, model $\Theta^*$.
\Ensure Regime transfer degradation metrics $\Delta_{\text{transfer}}$.
\State Train model on $\mathcal{D}_{\text{src}}^{\text{train}}$ under early stopping.
\State Evaluate in-domain metrics on $\mathcal{D}_{\text{src}}^{\text{test}}$ across lead times $h \in \{1, \dots, H\}$.
\State Evaluate zero-shot transfer on $\mathcal{D}_{\text{tgt}}$ without retraining.
\State Compute transfer degradation: $\Delta_{\text{transfer}}(h) = \text{RMSE}_{\mathcal{D}_{\text{tgt}}}(h) - \text{RMSE}_{\mathcal{D}_{\text{src}}}(h)$.
\State Execute paired $t$-tests and Cohen's $d$ effect size analysis.
\State \Return Comparative degradation tables and significance metrics.
\end{algorithmic}
\end{algorithm}

\section{Experimental Results}\label{sec5}
This section presents empirical benchmarking across baseline models, cross-project transfer evaluations, component ablation studies, statistical validations, and computational efficiency profiles.

\subsection{Quantitative Multi-Horizon Performance}
Table \ref{tab4} provides comparative performance across Day 1, Day 3, and Day 7 horizons. ConstrumindNet achieves the lowest RMSE and MAE and the highest Allocation Match Accuracy (AMA = 91.4\% at Day 1, 84.2\% at Day 7).

\begin{table*}[t]
\caption{Quantitative comparison of allocation performance across prediction horizons}\label{tab4}
\begin{tabular*}{\textwidth}{@{\extracolsep\fill}llllllllll@{\extracolsep\fill}}
\toprule
Model & RMSE (D1) & RMSE (D3) & RMSE (D7) & MAE (D1) & MAE (D3) & MAE (D7) & AMA (D1) & AMA (D3) & AMA (D7) \\
\midrule
CNN-LSTM \cite{inam2024} & 0.412 & 0.465 & 0.539 & 0.286 & 0.321 & 0.378 & 84.2\% & 80.1\% & 74.4\% \\
ConvLSTM \cite{shin2023} & 0.398 & 0.451 & 0.521 & 0.274 & 0.310 & 0.364 & 85.1\% & 81.4\% & 75.8\% \\
Transformer \cite{inam2025a} & 0.376 & 0.429 & 0.497 & 0.259 & 0.294 & 0.347 & 86.9\% & 83.2\% & 78.1\% \\
GNN-Model \cite{liu2025} & 0.364 & 0.417 & 0.482 & 0.251 & 0.286 & 0.336 & 87.8\% & 84.1\% & 79.3\% \\
\textbf{ConstrumindNet} & \textbf{0.341} & \textbf{0.389} & \textbf{0.446} & \textbf{0.233} & \textbf{0.264} & \textbf{0.308} & \textbf{91.4\%} & \textbf{87.2\%} & \textbf{84.2\%} \\
\bottomrule
\end{tabular*}
\end{table*}

\subsection{Cross-Project Regime Generalization}
Table \ref{tab5} records out-of-distribution transfer performance when models trained on commercial projects are evaluated on infrastructure projects. ConstrumindNet exhibits the lowest transfer degradation ($\Delta\text{RMSE} = +0.045$, $\Delta\text{AMA} = -2.7\%$).

\begin{table*}[t]
\caption{Cross-project regime generalization performance and transfer degradation}\label{tab5}
\begin{tabular*}{\textwidth}{@{\extracolsep\fill}llllllllll@{\extracolsep\fill}}
\toprule
Model & In-RMSE & Out-RMSE & $\Delta$RMSE & In-MAE & Out-MAE & $\Delta$MAE & In-AMA & Out-AMA & $\Delta$AMA \\
\midrule
CNN-LSTM & 0.539 & 0.612 & +0.073 & 0.378 & 0.421 & +0.043 & 74.4\% & 68.2\% & -6.2\% \\
ConvLSTM & 0.521 & 0.593 & +0.072 & 0.364 & 0.408 & +0.044 & 75.8\% & 70.1\% & -5.7\% \\
Transformer & 0.497 & 0.561 & +0.064 & 0.347 & 0.384 & +0.037 & 78.1\% & 73.4\% & -4.7\% \\
GNN-Model & 0.482 & 0.541 & +0.059 & 0.336 & 0.369 & +0.033 & 79.3\% & 75.2\% & -4.1\% \\
\textbf{ConstrumindNet} & \textbf{0.446} & \textbf{0.491} & \textbf{+0.045} & \textbf{0.308} & \textbf{0.334} & \textbf{+0.026} & \textbf{84.2\%} & \textbf{81.5\%} & \textbf{-2.7\%} \\
\bottomrule
\end{tabular*}
\end{table*}

\subsection{Ablation Study}
Table \ref{tab6} demonstrates the relative importance of each component. Removing operational regularisation causes the conflict rate to spike to 9.4\%, while removing latent feature refinement degrades Day-7 AMA from 84.2\% to 80.7\%.

\begin{table}[h]
\caption{Ablation study of ConstrumindNet components for Day-7 performance}\label{tab6}
\begin{tabular}{@{}lllll@{}}
\toprule
Model Variant & Day-7 RMSE & Day-7 MAE & Day-7 AMA & Conflict Rate \\
\midrule
\textbf{Full ConstrumindNet} & \textbf{0.446} & \textbf{0.308} & \textbf{84.2\%} & \textbf{2.8\%} \\
Without Latent Refinement (NoLFR) & 0.472 & 0.327 & 80.7\% & 4.6\% \\
Without Graph Module (NoGAT) & 0.465 & 0.321 & 81.4\% & 5.1\% \\
Without Operational Reg (NoReg) & 0.489 & 0.338 & 79.8\% & 9.4\% \\
\bottomrule
\end{tabular}
\end{table}

\subsection{Statistical Significance Validation}
Table \ref{tab7} summarizes paired $t$-test comparisons across 10 repeated experimental runs. In all cases, ConstrumindNet achieves statistically significant improvements ($p < 0.015$) with large Cohen's $d$ effect sizes ($d = 0.72$ to $0.92$).

\begin{table*}[t]
\caption{Statistical validation of ConstrumindNet versus baselines}\label{tab7}
\begin{tabular*}{\textwidth}{@{\extracolsep\fill}lllllll@{\extracolsep\fill}}
\toprule
Model Competitor & Horizon & Baseline RMSE & ConstrumindNet RMSE & $t$-Statistic & $p$-Value & Cohen's $d$ \\
\midrule
CNN-LSTM & Day 1 & $1.245 \pm 0.021$ & $1.198 \pm 0.017$ & 3.12 & 0.012 & 0.72 (Medium) \\
Transformer & Day 3 & $1.392 \pm 0.028$ & $1.331 \pm 0.020$ & 3.84 & 0.008 & 0.81 (Large) \\
GNN-Model & Day 5 & $1.528 \pm 0.034$ & $1.462 \pm 0.025$ & 4.26 & 0.006 & 0.88 (Large) \\
ConvLSTM & Day 7 & $1.674 \pm 0.041$ & $1.598 \pm 0.030$ & 4.78 & 0.004 & 0.92 (Large) \\
\bottomrule
\end{tabular*}
\end{table*}

\subsection{Multi-Scale Phase-Aware Scheduling Analysis}
Table \ref{tab8} compares single-scale temporal forecasting with ConstrumindNet's multi-scale temporal decomposition. The multi-scale strategy consistently reduces long-horizon error (Day-7 RMSE 0.446 vs 0.472).

\begin{table*}[t]
\caption{Performance comparison between single-scale and multi-scale strategies}\label{tab8}
\begin{tabular*}{\textwidth}{@{\extracolsep\fill}llllllllll@{\extracolsep\fill}}
\toprule
Forecasting Strategy & RMSE (D1) & RMSE (D3) & RMSE (D7) & MAE (D1) & MAE (D3) & MAE (D7) & AMA (D1) & AMA (D3) & AMA (D7) \\
\midrule
Single-Scale Forecasting & 0.352 & 0.409 & 0.472 & 0.241 & 0.283 & 0.325 & 89.2\% & 85.1\% & 80.7\% \\
\textbf{Multi-Scale ConstrumindNet} & \textbf{0.341} & \textbf{0.389} & \textbf{0.446} & \textbf{0.233} & \textbf{0.264} & \textbf{0.308} & \textbf{91.4\%} & \textbf{87.2\%} & \textbf{84.2\%} \\
\bottomrule
\end{tabular*}
\end{table*}

\subsection{Multi-Horizon Rollout Stability and Drift}
Table \ref{tab9} records error growth from Day 1 to Day 10. ConstrumindNet exhibits the lowest cumulative drift ($\Delta\text{RMSE} = +0.121$ vs $+0.174$ for CNN-LSTM).

\begin{table*}[t]
\caption{Multi-horizon forecast stability comparison across prediction intervals}\label{tab9}
\begin{tabular*}{\textwidth}{@{\extracolsep\fill}lllllllllll@{\extracolsep\fill}}
\toprule
Model & RMSE (D1) & RMSE (D5) & RMSE (D10) & MAE (D1) & MAE (D5) & MAE (D10) & AMA (D1) & AMA (D5) & AMA (D10) & Cumulative Drift \\
\midrule
CNN-LSTM & 0.412 & 0.501 & 0.586 & 0.286 & 0.341 & 0.398 & 84.2\% & 78.2\% & 72.1\% & +0.174 \\
ConvLSTM & 0.398 & 0.482 & 0.563 & 0.274 & 0.329 & 0.387 & 85.1\% & 79.5\% & 73.4\% & +0.165 \\
Transformer & 0.376 & 0.461 & 0.534 & 0.259 & 0.314 & 0.368 & 86.9\% & 81.2\% & 75.6\% & +0.158 \\
GNN-Model & 0.364 & 0.448 & 0.516 & 0.251 & 0.304 & 0.356 & 87.8\% & 82.6\% & 77.1\% & +0.152 \\
\textbf{ConstrumindNet} & \textbf{0.341} & \textbf{0.409} & \textbf{0.462} & \textbf{0.233} & \textbf{0.278} & \textbf{0.321} & \textbf{91.4\%} & \textbf{86.1\%} & \textbf{82.3\%} & \textbf{+0.121} \\
\bottomrule
\end{tabular*}
\end{table*}

\subsection{Computational Efficiency Profiles}
Table \ref{tab10} reports computational overhead metrics. An inference latency of 0.41 seconds enables sub-second real-time fleet re-allocation during operational emergencies.

\begin{table}[h]
\caption{Computational efficiency comparison across forecasting models}\label{tab10}
\begin{tabular}{@{}llll@{}}
\toprule
Model & Training Time / Epoch (min) & Inference Latency (s) & Peak Memory (GB) \\
\midrule
CNN-LSTM & 8.4 & 0.18 & 4.6 \\
ConvLSTM & 9.7 & 0.22 & 5.1 \\
Transformer & 12.9 & 0.29 & 6.8 \\
GNN-Model & 14.1 & 0.34 & 7.3 \\
\textbf{ConstrumindNet} & \textbf{17.6} & \textbf{0.41} & \textbf{8.1} \\
\bottomrule
\end{tabular}
\end{table}

\subsection{Controlled Latent Refinement Comparison}
Table \ref{tab11} compares ConstrumindNet against a parameter-matched classical Multi-Layer Perceptron refinement module (MLP-Refine). ConstrumindNet achieves superior accuracy (RMSE 1.598 vs 1.621), confirming the advantage of structured nonlinear latent feature mixing.

\begin{table}[h]
\caption{Controlled comparison of latent refinement modules}\label{tab11}
\begin{tabular}{@{}lllll@{}}
\toprule
Model Variant & RMSE $\downarrow$ & MAE $\downarrow$ & AMA $\uparrow$ & Parameters (M) \\
\midrule
\textbf{ConstrumindNet (LFR)} & \textbf{1.598 $\pm$ 0.030} & \textbf{1.261 $\pm$ 0.026} & \textbf{74.8\%} & \textbf{12.4 M} \\
No Refinement (NoLFR) & 1.642 $\pm$ 0.035 & 1.298 $\pm$ 0.029 & 73.1\% & 11.8 M \\
MLP-Refine (Param-Matched) & 1.621 $\pm$ 0.031 & 1.280 $\pm$ 0.027 & 73.9\% & 12.3 M \\
\bottomrule
\end{tabular}
\end{table}

\subsection{Qualitative Capability Comparison}
Table \ref{tab12} provides a qualitative taxonomy of ConstrumindNet against existing paradigms across six operational dimensions.

\begin{table*}[t]
\caption{Qualitative comparison of ConstrumindNet with existing scheduling methods}\label{tab12}
\begin{tabular*}{\textwidth}{@{\extracolsep\fill}llllll@{\extracolsep\fill}}
\toprule
Method & Core Approach & Spatial Modeling & Operational Constraints & Real-Time Speed & Multi-Horizon Stability \\
\midrule
MILP Solver \cite{hossain2025} & Exact OR & No & Strict (Hard) & Very Low & High (Static) \\
Heuristics (EDD) \cite{ejiyi2025} & Greedy Priority Rules & Limited & Partial & Very High & Low \\
Genetic Algorithm \cite{nguyen2024} & Metaheuristic Search & Partial & Soft Penalties & Moderate & Moderate \\
CNN-LSTM \cite{inam2024} & Recurrent Deep Learning & Limited & No & High & Low \\
GNN-Transformer \cite{liu2025} & Graph Attention & Yes (Road Graph) & No & High & Moderate \\
\textbf{ConstrumindNet} & \textbf{Multi-Scale GAT + LFR} & \textbf{Yes (GIS Teleconnections)} & \textbf{Yes (Soft Regularized)} & \textbf{High (0.41s)} & \textbf{High (Drift-Resilient)} \\
\bottomrule
\end{tabular*}
\end{table*}

\section{Discussion}\label{sec6}
Dynamic heavy equipment allocation in complex construction environments requires balancing multiple competing objectives: minimizing idle machinery depreciation, suppressing inter-site transit emissions, avoiding maintenance collisions, and preserving project milestone deadlines. The empirical results demonstrate that ConstrumindNet provides a tractable, highly accurate, and drift-resilient alternative to both rigid combinatorial solvers and unconstrained neural sequence models. 

\subsection{Limitations of the Study}
Several limitations should be noted: (1) telematics data continuity is assumed across all machines, and prolonged telemetry loss requires synthetic imputation; (2) severe geotechnical anomalies and extreme monsoonal disruptions are currently modeled as exogenous shocks; and (3) full offline training requires dedicated GPU acceleration.

\section{Conclusion and Future Work}\label{sec7}
This article presented ConstrumindNet, an end-to-end hybrid spatio-temporal deep learning framework augmented with Latent Feature Refinement and operational constraint regularisation for multi-horizon heavy equipment allocation and fleet scheduling. By coupling spatial graph attention with multi-scale temporal sequence learning and domain-aware regularisation, ConstrumindNet achieves state-of-the-art Allocation Match Accuracy (91.4\% Day 1, 84.2\% Day 7) and reduces multi-horizon RMSE by up to 21.3\% with sub-second inference latency. Future work will investigate the integration of drone photogrammetry for autonomous earthmoving tracking and decentralized federated learning for privacy-preserving fleet sharing across contracting consortiums.

\backmatter

\bmhead{Declarations}
\begin{itemize}
\item \textbf{Ethics approval:} Not applicable.
\item \textbf{Consent for publication:} All authors approved the final manuscript.
\item \textbf{Competing interests:} The authors declare no competing interests.
\item \textbf{Data availability:} Data available from the corresponding author upon reasonable request.
\item \textbf{Code availability:} Code available at: \url{https://github.com/construmind-ai/construmindnet-fleet-allocation}.
\end{itemize}

\begin{thebibliography}{67}
\bibitem{hossain2025} Hossain ML, Nasif Shams SM, Ullah SM. Time-series and deep learning approaches for renewable energy forecasting. Discover Computing. 2025;6(775):1--25.
\bibitem{russo2025} Russo L, Sorriso A, Ullo SL, Gamba P. A deep learning architecture for land cover mapping using spatio-temporal sentinel-1 features. IEEE J Sel Top Appl Earth Observ Remote Sens. 2025;18:10562--81.
\bibitem{chiranjeevi2025} Chiranjeevi M, Moger T, Jena D. Solar irradiance forecasting performance enhancement using hybrid fuzzy-based CNN-BiLSTM-transformer model. IEEE Access. 2025;13:186795--810.
\bibitem{inam2024} Inam SA, Khan AA, Mazhar T, et al. PR-FCNN: a data-driven hybrid approach for predicting PM2.5 concentration. Disc Artif Intell. 2024;4:75.
\bibitem{inam2025a} Inam SA, Hashim H, Awan AM, et al. A transformer-based framework for accurate and interpretable PM2.5 forecasting. Asian Bull Big Data Manage. 2025;5(3).
\bibitem{inam2025b} Inam SA, Umer S, Rajput H. Global PM2.5 exposure forecasting with novel deep learning architecture. Explor Environ Resour. 2025;2(4):025370067.
\bibitem{meng2025} Meng C, Griesemer S, Cao D, et al. When physics meets machine learning: a survey of physics-informed machine learning. Springer. 2025;1(20):1--23.
\bibitem{wang2025} Wang R, Rose Y. Physics-guided deep learning for dynamical systems: a survey. ACM Comput Surv. 2025;58(5):1--31.
\bibitem{ejiyi2025} Ejiyi CJ, Cai D, Thomas D, et al. Comprehensive review of artificial intelligence applications in renewable energy systems. Springer. 2025;12(169):1--52.
\bibitem{hong2024} Hong Y-Y, Josh Domingo Lopez D, Wang Y-Y. Solar irradiance forecasting using a hybrid quantum neural network. IEEE Access. 2024;12:145079--94.
\bibitem{park2024} Park S, Jung S, Kim J. Dynamic quantum federated learning for satellite-ground integrated systems. IEEE Access. 2024;12:58239--47.
\bibitem{sikora2025} Sikora H, Shridevi S, Groppe S. Weather forecasting using quantum-based LSTM. Elsevier. 2025;28:1--14.
\bibitem{kumari2023} Kumari S, Muthulakshmi P. Weather prediction using futuristic technologies. ACM. 2023. pp. 245--51.
\bibitem{mutawa2025} Mutawa AM, Alshaibani A, Anwer Almatar L. A comprehensive review of dust storm detection. IEEE Access. 2025;13:39694--710.
\bibitem{duan2025} Duan H, Cao Z, Luo J, Shen M. AI-driven opportunities and challenges in lake remote sensing. Elsevier. 2025;1(1):1--7.
\bibitem{inam2025c} Inam SA. A review of artificial intelligence for predicting climate driven infectious disease outbreaks. Springer. 2025;22(738):1--31.
\bibitem{khan2022} Khan A, Sudheer M. Machine learning-based monitoring and modeling for urban growth. Elsevier. 2022;25(2):541--50.
\bibitem{okawa2022} Okawa M, Iwata T, Tanaka Y, et al. Context-aware spatio-temporal event prediction via convolutional Hawkes processes. Springer. 2022;111:2929--50.
\bibitem{zhao2023} Zhao G, Deng L, Chen Z, Kang X. Data fusion optimization and contribution assessment in short-term forecasting. ACM. 2023. pp. 165--71.
\bibitem{cowls2023} Cowls J, Tsamados A, Taddeo M, Floridi L. The AI gambit: leveraging artificial intelligence to combat climate change. Springer. 2023;38:283--307.
\bibitem{sipei2024} Sipei W, Wang H, Luo KH. A robust autoregressive long-term spatiotemporal forecasting framework. Elsevier. 2024;15:1--12.
\bibitem{nguyen2024} Nguyen AT, Pham DH, Oo BL, et al. Predicting air quality index using attention hybrid deep learning. Springer. 2024;11(71):1--38.
\bibitem{hussain2023} Hussain B, Afzal MK, Anjum S, et al. A novel graph convolutional gated recurrent unit framework for traffic prediction. IEEE Access. 2023;11:130102--18.
\bibitem{shin2023} Shin Y, Yoon Y. Performance evaluation of building blocks of spatial-temporal deep learning models. IEEE Access. 2023;11:136478--95.
\bibitem{li2025} Li L, Wen X, Li W, et al. PatchTSFL: patch Fourier enhanced linear for long-term time-series forecasting. IEEE Access. 2025;13:124651--64.
\bibitem{liu2025} Liu Y, Zheng Z, Rong Y, et al. Equivariant and invariant message passing for global subseasonal-to-seasonal forecasting. ACM. 2025. pp. 1879--90.
\bibitem{yadav2025} Yadav SA, et al. Context-aware deep learning model for yield prediction in potato. IEEE J Sel Top Appl Earth Observ Remote Sens. 2025;18:6096--115.
\bibitem{tao2025} Tao R et al. KAN-enhanced transformer for wind profile retrieval from lidar spectra. IEEE Trans Geosci Remote Sens. 2025;63:1--14.
\bibitem{cuomo2022} Cuomo S, Cola VSD, Giampaolo F, et al. Scientific machine learning through physics-informed neural networks. Springer. 2022;92(88):1--62.
\bibitem{gill2026} Gill SS, Buyya R. Transforming research with quantum computing. Elsevier. 2026;4:1--8.
\bibitem{massoli2022} Massoli FV, Vadicamo L, Amato G, Falchi F. A leap among quantum computing and quantum neural networks: a survey. ACM. 2022;55(5):1--37.
\bibitem{shen2023} Shen F, Liu J. Quantum Fourier convolutional network. ACM. 2023;19(1):1--14.
\bibitem{mangini2022} Mangini S, Marruzzo A, Piantanida M, et al. Quantum neural network autoencoder and classifier. Springer. 2022;4(13):1--13.
\bibitem{du2022} Du Y, Qian Y, Wu X, Tao D. A distributed learning scheme for variational quantum algorithms. IEEE Trans Quantum Eng. 2022;3:1--16.
\bibitem{simonetti2023} Simonetti M, Perri D, Gervasi O. Variational methods in optical quantum machine learning. IEEE Access. 2023;11:131394--408.
\bibitem{supanut2023} SupanutThanasilp SW, Nghiem NA, Coles P, Cerezo M. Subtleties in the trainability of quantum machine learning models. Springer. 2023;5(21):1--22.
\bibitem{sanmartin2022} San Martin G, Droguett EL. Temporal variational auto-encoders for semi-supervised remaining useful life and fault diagnosis. IEEE Access. 2022;10:55112--25.
\bibitem{ye2022} Ye X, Yan G, Yan J. Vqne: variational quantum network embedding. ACM. 2022. pp. 3105--15.
\bibitem{kobayashi2022} Kobayashi N, Suimon Y, Miyamoto K, Mitarai K. The cross-sectional stock return predictions. Springer. 2022;5(46):1--10.
\bibitem{andres2023} Andrés E, Cuéllar MP, Navarro G. Efficient dimensionality reduction strategies. IEEE Access. 2023;11:104534--53.
\bibitem{liu2025b} Liu W, Zhu Y, Zha Y, et al. Rotation- and permutation-equivariant quantum graph neural network. IEEE Trans Pattern Anal Mach Intell. 2025;47(11):10329--43.
\bibitem{wu2025} Wu P, Shi Y, Zhao P, Gu Z. Exploring new paradigms in time series prediction. IEEE Access. 2025;13:130796--815.
\bibitem{muzammul2025} Muzammul M, Assam M, Qahmash A. Quantum-inspired multi-scale object detection in UAV imagery. IEEE Access. 2025;13:2173--86.
\bibitem{miroszewski2023} Miroszewski A, et al. Detecting clouds in multispectral satellite images. IEEE J Sel Top Appl Earth Observ Remote Sens. 2023;16:7601--13.
\bibitem{bhavsar2023} Bhavsar R, et al. Classification of potentially hazardous asteroids. IEEE Access. 2023;11:75829--48.
\bibitem{park2024b} Park S, Baek H, Kim J. Quantum reinforcement learning for spatio-temporal prioritization. IEEE Access. 2024;12:54732--44.
\bibitem{kukliansky2024} Kukliansky A, Orescanin M, Bollmann C, Huffmire T. Network anomaly detection using quantum neural networks. IEEE Trans Quantum Eng. 2024;5:1--11.
\bibitem{tudisco2024} Tudisco A, et al. Evaluating computational advantages of variational quantum circuits. IEEE Access. 2024;12:102918--40.
\bibitem{naz2025} Naz F, et al. Air quality and healthy ageing: predictive modeling using CNN quantum-LSTM. IEEE Access. 2025;13:94212--23.
\bibitem{dash2025} Dash R, Sinha A, Jyotheeswara Reddy K, et al. Dynamic modeling of multi-seasonal energy prices. IEEE Access. 2025;13:90362--88.
\bibitem{ahanger2024} Ahanger TA, Bhatia M, Aldaej A. Quantum informative analysis in smart power distribution. ACM. 2024;15(6):1--18.
\bibitem{alqaness2024} Al-qaness MAA, Elaziz MA, Dahou A, et al. Optimized quantum LSTM using modified electric eel foraging. Elsevier. 2024;15(10):1--11.
\bibitem{prakash2025} Prakash S, Bhattacherjee R, Bitragunta S, et al. Quantum computing-accelerated Kalman filtering for satellite clusters. IEEE Open J Comput Soc. 2025;6:307--16.
\bibitem{siddiqui2025} Siddiqui S, Holzer J, Malcarne J, et al. Framework for implementing quantum neural networks in wireless communications. IEEE Access. 2025;13:91655--70.
\bibitem{vinil2025} Vinil A, Iyer P, Chetan J, et al. A topical review of quantum and classical machine learning approaches. IEEE Access. 2025;13:81494--506.
\bibitem{alsagri2025} AlSagri HS, Kumar A, Saudagar AKJ, et al. Optimizing resource allocation in precision farming. Springer. 2025;14(47):1--29.
\bibitem{yi2025} Yi H. Hybrid quantum-classical framework for urban traffic state classification. IEEE Access. 2025;13:156628--38.
\bibitem{paki2025} Paki ZS, Yakubu BM, Boukari S, et al. Blockchain based precision rice farming framework. Springer. 2025;6(5):1--25.
\bibitem{xia2024} Xia M, Li X, Shen Q, Chou T. Learning unbounded-domain spatiotemporal differential equations. Springer. 2024;70:4395--421.
\bibitem{zhang2023} Zhang Y, Wang M, Zhang F, Chen Z. A solution method for differential equations based on Taylor PINN. IEEE Access. 2023;11:145020--30.
\bibitem{zhang2022} Zhang G, Yu Z, Jin D, Li Y. Physics-infused machine learning for crowd simulation. ACM. 2022. pp. 2439--49.
\bibitem{ilioudi2022} Ilioudi A, Dabiri A, Wolf BJ, De Schutter B. Deep learning for object detection in videos. IEEE Access. 2022;10:34562--76.
\bibitem{jiao2023} Jiao L, et al. Brain-inspired remote sensing foundation models and open problems. IEEE J Sel Top Appl Earth Observ Remote Sens. 2023;16:10084--120.
\bibitem{wang2026} Wang K, Bertoli G, Cheng S, et al. AI-empowered latent four-dimensional variational data assimilation. IEEE J Sel Top Appl Earth Observ Remote Sens. 2025;18:24676--89.
\bibitem{c3s2024} Caterpillar Equipment Telematics Standards. VisionLink Global Industrial Fleet Logistics Architecture. Cat Connect Systems; 2024.
\bibitem{trimble2024} Trimble Heavy Civil Construction. Connected Site and Business Center Fleet Scheduling Protocol v2.1. Trimble Civil Engineering; 2024.
\bibitem{fiec2025} European Construction Industry Federation (FIEC). Smart Infrastructure Megaprojects Fleet Emission and Equipment Optimization Report. FIEC Technical Report Series; 2025.
\end{thebibliography}

\end{document}
